% Propietario conservado: /Users/ruben/Documents/New project/output/REVISION_ARGUMENTAL_APERTURAS_CIERRES_20260915/VI/delivery/MANUSCRIPT_VI_ES_EN_COMPLETE/source_es/sections/10_sectores_topologicos.tex
\section{Persistance, structures de fusion et réalisations spectrales}
\label{vi:vi:sec:topologia}

L’atlas et les opérateurs de masse admettent différentes représentations. La
classification de ces représentations conserve la donnée qui caractérise
chacune : une règle de multiplication, une conjugaison réelle, une phase
d’échange ou une mesure spectrale. Cette section situe les secteurs
topologiques et les réalisations neutres dans la même architecture.
Les coordonnées de masse s’obtiennent une fois que la représentation et sa
dynamique sont déterminées.

La différence entre ces étapes est particulièrement visible dans le secteur
aréal--volumétrique. Une matrice d’incidence peut fixer un anneau de fusion
et sa dimension positive ; l’énergie d’une excitation dépend également de
l’opérateur qui agit dans la représentation. La première détermination conserve
un contenu exact lorsqu’on étudie plusieurs réalisations dynamiques
du même anneau. La distinction maintient l’information structurelle qui
précède la grandeur.

\subsection{L’anneau basé de l’incidence aréale--volumétrique}
\label{vi:vi:subsec:fusion-av}

L’incidence primitive des feuillets, obtenue à partir du calendrier du TPK, est
\[
 F_{\mathrm{av}}=\begin{pmatrix}0&1\\1&1\end{pmatrix}.
\]
Dans la base $(1,\tau)$, la première colonne exprime
$\tau\cdot1=\tau$ ; l’autre exprime $\tau\cdot\tau=1+\tau$.
L’incidence induit ainsi l’anneau basé
\[
 \mathcal F_{\mathrm{av}}=\mathbb Z[\tau]/(\tau^2-\tau-1),
 \qquad \mathcal F_{\mathrm{av}}^+
 =\{a+b\tau:a,b\in\mathbb Z_{\geq0}\}.
\]
Le cône positif enregistre des multiplicités. L’opération est interne aux
deux types et n’utilise pas une masse comme coefficient.

\begin{proposition}[Produit et dimension positive]
\label{vi:vi:prop:fusion-av}
On a
\[
 (a+b\tau)(c+d\tau)
 =(ac+bd)+(ad+bc+bd)\tau.
\]
Il existe un unique homomorphisme d’anneaux $d:\mathcal F_{\mathrm{av}}
\to\mathbb R$ avec $d(1)=1$ et $d(\tau)>0$. Sa valeur en $\tau$
est la coordonnée d’or reconnue dans la réalisation de $F_{\mathrm{av}}$.
\end{proposition}
\begin{proof}
Développer le produit et substituer $\tau^2=1+\tau$ donne les coefficients.
Un homomorphisme unital est déterminé par $x=d(\tau)$, avec
$x^2=x+1$. Ses racines sont $(1+\sqrt5)/2$ et $(1-\sqrt5)/2$ ;
seule la première est positive. L’évaluation en cette racine définit un
homomorphisme du quotient. L’équation identifie la dimension de cette
représentation ; la coordonnée régionale HMT a été construite auparavant.
\end{proof}

Avec $F_0=0$, $F_1=1$ et $F_{n+1}=F_n+F_{n-1}$, les puissances satisfont
\[
 \tau^n=F_{n-1}+F_n\tau\qquad(n\geq1).
\]
Le cas initial est immédiat ; multiplier par $\tau$ applique à nouveau
l’incidence et produit l’étape de récurrence. Les coefficients conservent
ainsi leur origine comme dénombrements de composition des feuillets.

Une représentation catégorique ajoute des espaces de morphismes et un
associateur à ces multiplicités. Une représentation tressée ajoute
des opérateurs d’échange compatibles. La réalisation matérielle ajoute
l’hamiltonien, ses domaines et la préparation qui sélectionne un
secteur spectral. Si elle possède un écart $\Delta E>0$ et une vitesse effective
$c_{\mathrm{eff}}$, sa coordonnée de masse d’écart est
\[
 m_{\mathrm{brecha}}=\Delta E\,c_{\mathrm{eff}}^{-2}.
\]
La dimension positive de l’anneau et l’écart sont des sorties d’opérateurs
différents. La règle de fusion conserve son contenu sans fixer, à
elle seule, une masse universelle pour toutes ses réalisations.

\subsection{Dimension et type de représentation}

Dans une représentation pointée, tous les objets simples sont
inversibles. Si $X$ est un simple inversible et $d$ une dimension
positive compatible avec les duaux, alors
\[
 1=d(X\otimes X^*)=d(X)d(X^*)=d(X)^2,
\]
et $d(X)=1$. La dimension fournit donc un contrôle des
changements de représentation.

\begin{proposition}[Obstruction de dimension dans le secteur pointé]
\label{vi:vi:prop:apuntada}
Un simple de dimension $\varphi$ ou $\sqrt2$ ne peut pas être identifié,
en préservant la dimension, à un simple d’une représentation pointée.
En particulier, les $81$ simples inversibles du centre non tordu
basé sur $\mathbb Z/9\mathbb Z$ ne fournissent pas cette identification.
\end{proposition}
\begin{proof}
Les simples inversibles ont pour dimension $1$, distincte des deux
dimensions. Dans le centre non tordu abélien indiqué, un simple est
étiqueté par un élément du groupe et un caractère de son dual. Les deux
ensembles ont $9$ éléments. Les $81$ étiquettes sont inversibles
par la multiplication du groupe et du caractère.
\end{proof}

Ce recensement $81$ a un domaine distinct de l’atlas de $81$ cellules.
Une extension ou un foncteur vers un autre codomaine déclare les actions
et les applications qu’il transporte. L’incidence aréale--volumétrique conserve
son anneau exact et la réalisation topologique conserve les conditions
de son propre type.

\subsection{Autoconjugaison neutre et générateurs réels de CAR}
\label{vi:vi:subsec:majorana-categorias}

Une structure réelle sur un espace complexe est un opérateur antiunitaire
$\mathcal C$ avec $\mathcal C^2=I$. Sa partie fixe est un espace réel.
La commutation de l’opérateur de masse avec $\mathcal C$ préserve cette
partie réelle dans son domaine. La restriction d’une évolution à la partie
fixe s’exprime par la commutation de son groupe avec $\mathcal C$ ;
dans une évolution de Schr\"odinger, il faut distinguer cette condition
de la commutation de l’hamiltonien. Les deux conditions appartiennent à la
représentation matérielle de la section~\ref{vi:vi:subsec:antisecciones}.

\begin{proposition}[Conjugaison antiunitaire et évolution temporelle]
\label{vi:vi:prop:conjugacion-evolucion}
Soient $H$ autoadjoint, $\hbar>0$ et $\mathcal C$ une conjugaison
antiunitaire qui conserve $\mathcal D(H)$. Pour
$U(t)=\exp(-itH/\hbar)$, si
$\mathcal C H\mathcal C^{-1}=\epsilon H$, avec
$\epsilon\in\{1,-1\}$, on a
\begin{equation}
 \mathcal C U(t)\mathcal C^{-1}=U(-\epsilon t).
 \label{vi:vi:eq:conjugacion-grupo-temporal}
\end{equation}
Le groupe $U(t)$ préserve la partie fixe de $\mathcal C$ pour tout $t$
si et seulement s’il commute avec $\mathcal C$, de manière équivalente si
$\mathcal C H\mathcal C^{-1}=-H$.
Si, en revanche, $\mathcal C H=H\mathcal C$, un vecteur fixe sous
$\mathcal C$ reste fixe sous celle-ci pendant toute l’évolution
exactement lorsqu’il appartient à $\ker H$.
\end{proposition}
\begin{proof}
La conjugaison antiunitaire envoie $i$ sur $-i$. Le calcul fonctionnel
spectral fournit
\[
 \mathcal C\exp(-itH/\hbar)\mathcal C^{-1}
 =\exp\bigl(it\mathcal C H\mathcal C^{-1}/\hbar\bigr),
\]
ce qui prouve \eqref{vi:vi:eq:conjugacion-grupo-temporal}. Si $U(t)$ commute
avec $\mathcal C$, il transporte ses vecteurs fixes en vecteurs fixes.
Réciproquement, tout vecteur complexe s’écrit de manière unique comme
$x+iy$ avec $x,y$ fixes sous $\mathcal C$. La préservation de cette partie
réelle implique, par linéarité de $U(t)$ et antilinéarité de
$\mathcal C$, que les deux commutent. L’unicité du générateur
autoadjoint du groupe unitaire identifie alors
$\mathcal C H\mathcal C^{-1}$ à $-H$.

Dans le cas commutatif $\epsilon=1$, pour $\mathcal C\psi=\psi$,
on a $\mathcal C U(t)\psi=U(-t)\psi$. Que $U(t)\psi$ soit fixe
pour tout $t$ équivaut à $U(2t)\psi=\psi$ pour tout $t$. Par le
calcul spectral, cela équivaut à ce que la mesure spectrale de $\psi$
soit supportée dans $\{0\}$, c’est-à-dire à $\psi\in\ker H$.
\end{proof}

Ainsi, la symétrie antiunitaire qui commute avec $H$ inverse le sens
temporel. La réalisation par des états réels conservés utilise
une autre condition sur le générateur. Cette distinction préserve la portée
de l’autoconjugaison algébrique et permet de spécifier séparément la
dynamique d’un champ de Majorana.

\begin{proposition}[Autoconjugaison et charge]
\label{vi:vi:prop:carga-real}
Soient $Q$ un opérateur de charge autoadjoint et $\mathcal C$ une
conjugaison qui transporte son domaine et satisfait
$\mathcal C Q\mathcal C^{-1}=-Q$. Si $\psi\ne0$,
$Q\psi=q\psi$ et $\mathcal C\psi=\psi$, alors $q=0$.
\end{proposition}
\begin{proof}
L’autoadjonction implique $q\in\mathbb R$. Appliquer $\mathcal C$
à l’équation propre donne $-Q\psi=q\psi$. Avec l’équation initiale,
on obtient $2q\psi=0$, donc $q=0$.
\end{proof}

La neutralité est nécessaire dans ce problème d’autoconjugaison ;
le noyau de $Q$ ne spécifie pas, à lui seul, un opérateur antiunitaire
ni sa compatibilité avec le groupe d’évolution. Le critère s’applique aux
blocs neutriniques après la construction de leurs opérateurs et de leur mélange.
Le mot dodécaphasique fixe, le spineur réel et la représentation de
charge conservent ainsi leurs domaines respectifs.

La complétion fermionique fournit une autre réalisation. Dans
l’espace extérieur des modes déclarés, $a_j^*$ est le produit
extérieur par le vecteur de base et $a_j$ son adjoint. L’antisymétrie
extérieure donne
\[
 \{a_j,a_k\}=0,\qquad \{a_j,a_k^*\}=\delta_{jk}I.
\]

\begin{proposition}[Générateurs autoadjoints de CAR]
\label{vi:vi:prop:car-real}
Les opérateurs
\[
 \gamma_{2j-1}=a_j+a_j^*,\qquad
 \gamma_{2j}=-i(a_j-a_j^*)
\]
sont autoadjoints et vérifient $\{\gamma_r,\gamma_s\}=2\delta_{rs}I$.
\end{proposition}
\begin{proof}
L’adjonction laisse fixe chaque opérateur. Pour le même mode,
$a_j^2=(a_j^*)^2=0$ et $a_ja_j^*+a_j^*a_j=I$, donc chaque
carré vaut $I$. L’anticommutateur entre les deux combinaisons
du même mode s’annule par développement. Pour des modes distincts,
tous les anticommutateurs élémentaires s’annulent.
\end{proof}

Ce sont des générateurs réels d’une structure de Clifford. Un mode
zéro topologique nécessite en outre que la dynamique réalise une excitation
localisée d’énergie nulle, séparée du continu par un écart et
stable sous les perturbations considérées. La propriété algébrique
et la propriété spectrale sont reliées par la réalisation dynamique.

\subsection{Phases d’échange et enroulement d’ordre six}

Pour \(n\ge2\), dans le lecteur d’ordre six, la somme des exposants définit
$w_6:B_n\to\mathbb Z/6\mathbb Z$. Les relations d’Artin
conservent cette somme : la commutation des générateurs distants maintient
deux termes et la relation à trois générateurs en conserve trois.

\begin{proposition}[Caractères du lecteur d’enroulement]
\label{vi:vi:prop:devanado-seis}
Les caractères unitaires du quotient sont
\[
 \chi_k(b)=\exp\!\left(\frac{2\pi i k}{6}w_6(b)\right),
 \qquad k=0,1,\ldots,5.
\]
Une représentation scalaire de cette forme passe au groupe des
permutations exactement pour $k=0$ ou $k=3$.
\end{proposition}
\begin{proof}
Un caractère cyclique est déterminé par l’image $q$ de $1$,
avec $q^6=1$. Pour passer aux permutations, il doit satisfaire la
relation supplémentaire $\sigma_j^2=1$, qui exige $q^2=1$.
Parmi les racines sixièmes, seules $q=1,-1$ vérifient cette condition.
Toutes deux satisfont toutes les relations du groupe des permutations.
\end{proof}

Les quatre autres caractères conservent une information d’enroulement.
Une espèce matérielle nécessite que le lecteur se réalise sur une
multisection et que la dynamique conserve sa structure. L’inversion
de route remplace $w_6$ par $-w_6$ et transporte chaque caractère en son
conjugué complexe.

\subsection{Stabilité et visibilité d’une section}

La positivité démontrée dans la
section~\ref{vi:vi:sec:operador-masa} a une conséquence spectrale.
Pour un opérateur autoadjoint positif $M$, le calcul fonctionnel donne
$\operatorname{Spec}(M^2)\subset[0,\infty)$. Dans une réalisation qui
identifie $M^2$ à l’observable de masse au carré sur un
espace physique à produit positif, une section propre persistante
ne peut pas avoir de valeur propre négative de cette observable. Une direction
instable de la linéarisation d’un fond appartient à un autre problème
spectral, qui est étudié avec son domaine et sa dynamique.

La visibilité dépend du couplage. Si $P_X$ est le support d’une
section et $P_a$ le projecteur d’un canal, $P_aP_X=0$ exprime
l’orthogonalité des supports. Le produit non nul exprime une incidence
possible ; son intensité est déterminée par l’opérateur de couplage.
Deux projecteurs peuvent commuter et avoir un produit non nul, de sorte
que le commutateur ne remplace pas cette incidence.

Dans un bloc neutre, une section stérile par rapport au canal faible
appartient au noyau de son couplage. Si ce noyau est invariant
pour l'opérateur de masse, la restriction possède une mesure spectrale
propre. La profondeur de l'atlas, la neutralité et l'absence de couplage
faible sont des propriétés distinctes vérifiées sur la même
section. Les possibilités de représentation sont conservées sans
transformer une étiquette de famille en masse externe.

L’organisation maintient ainsi l’ordre de l’argument : l’incidence
détermine les multiplicités ; les relèvements déterminent
les conjugaisons et les phases ; la représentation et l’opérateur matériel
déterminent la sortie spectrale. La confrontation s’applique ensuite à
cette sortie, avec ses unités et son couplage déclarés.
